\chapter{Derivational Morphology}\label{ch:derivational}


\section{Nominal derivation}

\subsection{Diminutives and augmentatives}


\subsection{Agent nouns: \ird{-(h)ár} and \ird{-ist}}

Agent nouns in \ird{-(h)ár} and \ird{-ist} denote the habitual or
professional performer of an action; both suffixes are productive.

% [Examples: fisherman, teacher, activist, etc.]
% Note: the agent nominal -kou (from the antipassive stem) is a verbal
% derivation documented in §\,\ref{sec:nominalised}; the -ár/-ist forms
% are purely nominal derivation.

\subsection{Place nouns: \ird{-(i)mašt}}


\subsection{Abstract nouns: \ird{-(i)žnám}}


\section{Verbal derivation}

Unlike the valency operations described in Chapter~\ref{ch:verbs}, the
processes treated here do not occupy obligatory positions in the finite verbal
template. They form specialised lexical stems or stem variants, which then
take the ordinary inflectional morphology of the language. In the modern
language, denominal verbalisation and a restricted class of lexical preverbs
form new verbal stems.

\subsection{Denominal verbalisation}

Among overt noun-to-verb strategies, suffixation in \ird{-ěvat-} and
\ird{-it-} is the most productive. These are the predominant verbalising
formatives with learned, foreign, and recently borrowed nominal bases.
When a foreign stem is adapted directly as a verb, \ird{-ěvat-} is used
exclusively. This construction selects the palatal grade at the nominal
stem boundary; its output is written with explicit \ird{ť} or \ird{ň}
where those consonants result (§\,\ref{sec:morphophon}). The grade selected
by \ird{-it-} remains unspecified pending paradigmatic evidence: its initial
\ird{i} alone does not establish a categorical mutation rule.

\pex
\a \ird{kontakt} $\rightarrow$ \ird{kontakťevat-} \trsl{contact; get in touch}
\a \ird{telefon} $\rightarrow$ \ird{telefoňevat-} \trsl{telephone}
\xe

\subsubsection{Relational statives in \ird{-(e)zn-}}

A semantically specialised but productive denominal pattern adds
\ird{-(e)zn-} to a nominal base to form an intransitive relational stative
verb. The core meaning is \trsl{be X-like; belong to the domain of X; have
the character, style, material, or social colouring of X}. These verbs
conjugate exactly like the ordinary stative class described in
§\,\ref{sec:stative}, and they likewise admit the attributive derivative in
\ird{-á} and the manner adverb in \ird{-ě}
(§§\,\ref{sec:stative-attr}--\ref{sec:stative-adv}). Aspect follows the
usual stative pattern: the continuous gives plain state predication, the
progressive and perfective usually describe entry into the state, and the
retrospective describes a former state. The construction selects the restricted
velar-front grade: stem-final \ird{k} surfaces as \ird{c}, as in
\ird{barok} $\rightarrow$ \ird{barocezna}, while coronals remain plain,
as in \ird{kmětezna} and \ird{svícezna}. The independently added manner
suffix \ird{-ě} selects the palatal grade at the final consonant of the
derived stative stem, yielding \ird{-zňe}.

Table~\ref{tab:relstatives} lists representative continuous forms together
with their derived attributives and manner adverbs.

\begin{table}[ht]
  \footnotesize
  \caption{Relational stative verbs in \ird{-(e)zn-}.}\label{tab:relstatives}
  \begin{tabularx}{\textwidth}{L l l l L}
    \toprule
    {\sc base noun} & {\sc relational stative verb} & {\sc attributive} & {\sc manner adverb} & {\sc core sense} \\
    \midrule
    \ird{tvír} \trsl{court, courtly sphere} & \ird{tvírzna} & \ird{tvírzná} & \ird{tvírzňe} & \trsl{be courtly; belong to the courtly register} \\
    \ird{ščech} \trsl{guild, craft corporation} & \ird{ščehezna} & \ird{ščehezná} & \ird{ščehezňe} & \trsl{be guild-like; corporate, artisanal, guild-marked} \\
    \ird{mněra} \trsl{faith, creed} & \ird{mněrazna} & \ird{mněrazná} & \ird{mněrazňe} & \trsl{be confessional; marked by a creed or religious orientation} \\
    \ird{kaměn} \trsl{stone} & \ird{kaměnzna} & \ird{kaměnzná} & \ird{kaměnzňe} & \trsl{be stony; have a stone-like character; be stone-styled} \\
    \ird{sylo} \trsl{cloth, textile} & \ird{sylozna} & \ird{sylozná} & \ird{sylozňe} & \trsl{be textile-like; be clothed, draped, or fabric-marked} \\
    \ird{barok} \trsl{baroque; baroque style} & \ird{barocezna} & \ird{barocezná} & \ird{barocezňe} & \trsl{be baroque in style; belong to a baroque aesthetic} \\
    \ird{zmok} \trsl{temple, shrine} & \ird{zmocezna} & \ird{zmocezná} & \ird{zmocezňe} & \trsl{be temple-like; be cultic, sacral, or shrine-associated} \\
    \ird{svíc} \trsl{peasant, smallholder, rustic estate} & \ird{svícezna} & \ird{svícezná} & \ird{svícezňe} & \trsl{be peasantlike; be rustic or village-class marked} \\
    \ird{kmět} \trsl{law, order, code} & \ird{kmětezna} & \ird{kmětezná} & \ird{kmětezňe} & \trsl{be juridical, code-bound, institutionalised} \\
    \bottomrule
  \end{tabularx}
\end{table}

\pex
\a \ird{tvírzná rěč} \trsl{courtly speech; speech of courtly register}
\a \ird{kaměnzná stěna} \trsl{a wall with a stony character} \lingcontext{a literal \trsl{stone wall} is expressed analytically, not by \ird{kaměnzná}}
\a \ird{barocezná dvorba} \trsl{a baroque-style building}
\xe

\subsection{Lexical preverbs}\label{sec:lexical-preverbs}

The bound preverbs \ird{de-}, \ird{mně-}, and \ird{pře-} form lexical
stems denoting directed motion, traversal, and replacement or revision,
respectively. Their productivity depends on the event structure of the base:
\ird{de-} combines with predicates admitting a goal, \ird{mně-} with
activities admitting an ordered path or domain, and the more restricted
\ird{pře-} chiefly with cognition, assessment, and transformation. A fourth
formative, privative \ird{ne-}, is confined to individually specified lexical
derivatives. Table~\ref{tab:lexical-preverbs} summarises these distributions;
the hyphenated forms are bound lexical stems, not finite verbs.

\begin{table}[htbp]
  \footnotesize
  \caption{Lexical preverbs and their principal distributions.}\label{tab:lexical-preverbs}
  \begin{tabularx}{\textwidth}{l L L}
    \toprule
    {\sc preverb} & {\sc contribution and distribution} & {\sc derivative} \\
    \midrule
    \ird{de-} & Goal-directed or endpoint-oriented event; productive with compatible motion predicates
      & \ird{nočk-} \trsl{run} $\rightarrow$ \ird{denočk-} \trsl{run to a destination} \\
    \ird{mně-} & Traversal of a path or ordered domain; productive where stages can be construed as a coherent course
      & \ird{ktař-} \trsl{read} $\rightarrow$ \ird{mněktař-} \trsl{read through}; \ird{vad-} \trsl{think} $\rightarrow$ \ird{mněvad-} \trsl{think through} \\
    \ird{pře-} & Replacement of an earlier outcome or revision of a judgement; restricted productivity
      & \ird{vad-} \trsl{think} $\rightarrow$ \ird{převad-} \trsl{reconsider; revise a judgement} \\
    \ird{ne-} & Privation, removal, or dispossession; lexically restricted
      & Meaning and selection specified for each derivative; no general \trsl{undo} or negative rule \\
    \bottomrule
  \end{tabularx}
\end{table}

In \ird{mněktař-}, the text supplies an ordered domain through which reading
proceeds, while \ird{mněvad-} extends the same organisation to the stages of
considering a problem. The availability of such a domain restricts the pattern:
\ird{mně-} cannot be added indiscriminately to every activity verb. The
revision sense of \ird{převad-} concerns an earlier judgement or outcome.
Privative \ird{ne-} likewise contributes a lexical
meaning of removal or deprivation wherever a derivative licenses it; it does
not express the failure or non-occurrence of the base event.

\subsubsection{Stem formation, aspect, and argument structure}
\label{sec:preverb-aspect}

A preverb and its base form one lexical constituent, to which the ordinary
verbal morphology of §\,\ref{sec:morph-template} subsequently attaches. The
schematic organisation is therefore:

\begin{center}
\small
\begin{tabular}{c}
(\Neg{}) + [preverb + root]$_{\text{stem}}$ + (Voice/Appl) + (Caus)\\
+ Asp.Align + (Mood) + (Evid)
\end{tabular}
\end{center}

The bracketed constituent occupies the stem position of the finite template;
prefixation supplies no additional inflectional category. In the analysis
adopted here, the lexical constituent is interpreted within the predicate,
below the functional heads which license participants, causation, and viewpoint
aspect. The suffixal hierarchy described in Chapter~\ref{ch:verbs} consequently
applies to prefixed and unprefixed stems alike.

The distinction between lexical event structure and grammatical viewpoint has
a close comparative parallel in Hungarian, where a telic predicate can receive
a progressive interpretation \cite[816--817]{kardosfarkas2022}. Hungarian
separable particles also alternate in position under focus, negation, and
progressive aspect \cite[815]{kardosfarkas2022}; Iridian preverbs remain
bound throughout these environments. The comparison concerns the independence
of the two aspectual dimensions, without extending the Hungarian syntactic
analysis to Iridian.

The aspectual paradigm of \ird{mněktař-} illustrates the independence of
lexical traversal from grammatical viewpoint. Reading through a text has a
potential endpoint, but neither the continuous nor the progressive asserts
that the endpoint has been reached. The contemplative projects the reading
event, whereas the perfective presents it as a bounded whole and the
retrospective relates the prior event to a subsequent vantage point. The
transitive forms in Table~\ref{tab:preverb-aspect} retain the text as patient;
the perfective and retrospective therefore include overt Voice \ird{-n-}.

\begin{table}[htbp]
  \footnotesize
  \caption{Transitive \ird{mněktař-} across the five aspects.}\label{tab:preverb-aspect}
  \begin{tabularx}{\textwidth}{l l L}
    \toprule
    {\sc aspect} & {\sc finite form} & {\sc interpretation} \\
    \midrule
    Continuous & \ird{mněktařa} & Be engaged in reading through a text; completion is not asserted \\
    Progressive & \ird{mněktařime} & Be actively proceeding through a text \\
    Contemplative & \ird{mněktařách} & Intend or expect to read through a text \\
    Perfective & \ird{mněktařnik} & Read through a text as a bounded completed event \\
    Retrospective & \ird{mněktařnac} & Have read through a text, viewed from a subsequent reference point \\
    \bottomrule
  \end{tabularx}
\end{table}

Prefixation itself neither licenses an additional argument nor changes case
assignment. Thus \ird{denočk-} retains the intransitive motion frame of
\ird{nočk-}, with any overt destination expressed by the ordinary goal
construction, while \ird{mněktař-} retains the transitive frame of
\ird{ktař-}. In the perfective and retrospective domains, an ordinary
transitive derivative takes Voice \ird{-n-} and ergative alignment; its
antipassive instead follows the independently established patient-suppressing
construction. The continuous, progressive, and contemplative retain nominative
alignment (§\,\ref{sec:vm-asp-case}). An exceptional lexical case frame or
stem alternant requires specification with the particular derivative; it does
not follow from the presence of a preverb.

The unprefixed \ird{ktař-} can itself occur in the perfective, which bounds
a reading event without adding the traversal meaning of \ird{mně-}.
Conversely, the continuous and progressive already provide internal viewpoints
on \ird{mněktař-}; prefixed stems require no separate imperfectivising
derivation. The phasal particles \ird{v} and \ird{s} retain their ordinary
scope over the predicate and contribute the interpretations described in
§\,\ref{sec:class1-particles}, independently of the lexical preverb.

\subsubsection{Bound status and lexical restrictions}

A lexical preverb remains adjacent to its base under focus, negation, and
changes of mood. It is written together with the inflected stem, and no nominal,
particle, or other syntactic constituent can intervene at the preverb--root
boundary. Indicative \ird{za-} attaches outside the derived stem, as in
\ird{zamněktařime} \trsl{is not reading through}; the lexical privative
\ird{ne-} occupies a different position and supplies no substitute for
indicative negation (§\,\ref{sec:vm-neg}). Productive derivation admits one
lexical preverb per base. An apparently complex borrowed or lexicalised stem,
such as \ird{prit-} or \ird{podohlet-}, is not thereby segmented into a
productive native prefix and root.

The preverb belongs to the minimal lexical prosodic word and receives initial
primary stress: \ird{mněktařime} and \ird{zamněktařime} are both stressed
on \ird{mně-}, while \ird{převadime} is stressed on \ird{pře-}
(§\,\ref{sec:stress}). Prefixation introduces no independent mutation or
quantity grade. The derivative follows the ordinary construction-selected
stem grades, suffixal consonant mutations, and cluster phonology, with any
lexically exceptional allomorph recorded separately. Each derived stem has its
own lexical entry, specifying its base, meaning, argument structure, and any
special alternants. The synchronic restriction to particular semantic classes
does not establish a history of origin from modern postpositions, nor does
surface resemblance to a borrowed stem establish productive segmentation.

\section{Compounding}\label{sec:compound-prosody}

A compound integrated as one lexical prosodic word receives primary stress on
its first syllable, after the stem grades and syllable structure of the complete
form have been established (§\,\ref{sec:stress}). Where the components
remain separate prosodic words, each retains its own lexical stress; phrasal
prominence may favour one component. A separately specified pattern of secondary
stress is not required by the initial-stress rule and is left open here pending
lexical evidence.

\section{Borrowing and assimilation patterns}\label{sec:loan-prosody}

A borrowing integrated into Iridian phonology follows the ordinary initial-stress
rule, independently of the stress and vowel quantity of its source. Source-language
stress may occur in deliberate foreign pronunciation of a name or quotation;
such a pronunciation is distinct from an integrated lexical exception. No general
class of exceptional native stress is posited here. Any loan that retains an
exceptional stress pattern in ordinary Iridian requires an explicit lexical
specification, together with its vowel quantities and syllabification.
